\section{Methods}

The empirical strategy is a property-level income--cost and scenario model. It links EPC-derived dwelling characteristics, synthetic LSOA income distributions, required energy costs, price-shock scenarios, domestic net zero intervention packages, and spatial-distributional outputs. The unit of analysis is property $p$ in LSOA $l$ and MSOA $m$. Where observed household incomes are unavailable, the model estimates the probability that a household drawn from the local residual-income distribution would face a given required-energy burden. This design is deliberately probabilistic: it does not claim to observe the income of the current occupant of each property.

The equations are chosen to match the conceptual problem rather than to maximise statistical complexity. The paper needs a model that can do four things transparently: estimate required energy affordability under current conditions; expose the same property-income combinations to future fuel-price shocks; distinguish average burden from upper-tail burden risk; and compare intervention packages only after technical, tenure, tariff, and governance constraints are represented. A property-level scenario model is appropriate because the key mechanisms operate at different scales: dwelling fabric and heating systems are property-level, income distributions are small-area approximations, tariffs and fuel prices are market-level, and feasibility constraints are mediated by tenure, leasehold, heritage, and local delivery. The model is therefore not a prediction of exact bills for named households. It is a structured counterfactual framework for asking which property--income--governance combinations are most exposed and which interventions reduce that exposure.

This approach is also deliberately modular. The income--cost engine establishes the baseline denominator and required-cost numerator. The price-shock engine changes unit prices, standing charges, and mitigation while holding the logic of required demand explicit. The intervention engine then changes demand, fuel structure, tariff access, and feasibility. Because each layer modifies a different part of the same burden equation, the resulting de-risking estimates remain interpretable: a lower risk score can be traced to lower required demand, lower unit-price exposure, improved tariff access, policy support, or greater feasibility. This is why the equations are sufficient for the paper's purpose even though they do not observe all household behaviours directly.

\subsection{Income--Cost Engine}

The income--cost engine answers the first empirical problem: how can current required-burden exposure be estimated when the dataset contains properties rather than verified household incomes? A single average income would suppress Westminster's local lower tail, while assigning one synthetic income to each property would imply false precision. The constrained lognormal specification provides a middle route. It keeps MSOA after-housing-cost income as an anchor, allows LSOA-level lower-tail differentiation through deprivation information, and returns probabilities rather than definitive household labels.

The implemented income layer estimates one constrained lognormal after-housing-cost income distribution for each LSOA:
\begin{equation}
Y_l \sim \operatorname{Lognormal}(\mu_l,\sigma_l^2),
\end{equation}
where $Y_l$ denotes annual residual household income. The distribution is parameterised by its mean $\bar{Y}_l$:
\begin{equation}
\mu_l=\ln(\bar{Y}_l)-\frac{1}{2}\sigma_l^2 .
\end{equation}
For each MSOA, the LSOA parameters are jointly estimated so that the property-weighted mean remains anchored to the MSOA after-housing-cost income, while the lower tail is differentiated by the LSOA income-deprivation signal. The LSOA income-deprivation term is treated as a robust soft constraint on the lower tail, not as an equality target. For a given regularisation parameter $\lambda$, the current objective is:
\begin{equation}
\begin{split}
\mathcal{L}_m(\lambda)=&
\left[
\ln \left(\frac{\sum_{l \in m} N_l \bar{Y}_l}{\sum_{l \in m} N_l}\right)
-\ln(\bar{Y}_m)
\right]^2 \\
&+\lambda \sum_{l \in m} w_l
\rho_{\eta}\left(F_l(z)-d_l\right),
\end{split}
\end{equation}
where $N_l$ is the number of EPC property records in LSOA $l$, $w_l=N_l/\sum_{l \in m}N_l$, $\bar{Y}_m$ is MSOA residual income, $z=19{,}009$ is the income threshold used in the current module, $F_l(\cdot)$ is the LSOA lognormal cumulative distribution function, and $d_l$ is the LSOA income-deprivation score. The robust loss is the pseudo-Huber L1 form:
\begin{equation}
\rho_{\eta}(d)=\sqrt{d^2+\eta^2}-\eta,
\qquad \eta=10^{-3}.
\end{equation}
The fixed $\eta$ is a numerical differentiability regulariser for optimisation rather than a substantive research parameter. The current implementation constrains $\sigma_l \in [0.10,1.40]$ and searches a log-spaced grid of 100 $\lambda$ values from $10^{-4}$ to $10^{-1}$. Official LSOA fuel-poverty rates are used only as an external diagnostic to select $\lambda$ by minimising property-weighted RMSE, with weighted MAE and then the larger $\lambda$ as tie-breakers; they are not entered directly as the optimisation target.

The May 2026 sensitivity run also varies the MSOA after-housing-cost income anchor across central, upper-confidence-limit, and lower-confidence-limit scenarios. These scenarios bracket income-anchor uncertainty for the baseline burden model. They are not income forecasts, and they do not replace the canonical property-level master until a deliberate rerun is approved.

Required energy cost is derived from SAP and floor area rather than from observed consumption, because actual expenditure may be suppressed by underheating or self-rationing. For property $p$, the SAP energy cost factor is:
\begin{equation}
\operatorname{ECF}_p=
\begin{cases}
(100-\operatorname{SAP}_p)/16.21, & \operatorname{SAP}_p \geq 43.265,\\
10^{(108.8-\operatorname{SAP}_p)/120.5}, & \operatorname{SAP}_p < 43.265.
\end{cases}
\end{equation}
The baseline required annual energy cost is then:
\begin{equation}
C_{p,0}=\operatorname{ECF}_p \frac{A_p+45}{0.2},
\end{equation}
where $A_p$ is total floor area. Baseline required energy burden for a household with income $Y_l$ is:
\begin{equation}
B_{p,0}(Y_l)=\frac{C_{p,0}}{Y_l}.
\end{equation}
The current module therefore estimates burden-threshold probabilities:
\begin{equation}
P_{p,0}(\tau)=\Pr \left(B_{p,0}(Y_l)>\tau \right)
=F_l\left(\frac{C_{p,0}}{\tau}\right).
\end{equation}
The implemented baseline output uses $\tau=0.10$. For modelled LSOA fuel-poverty comparison, properties with SAP grades A, B, or C are assigned a property poverty score of zero; other properties use $P_{p,0}(0.10)$. This score is a baseline affordability input, not a complete measure of household fuel price risk.

The threshold probability is useful because it reports exposure to a policy-relevant affordability boundary without claiming that the boundary is the whole concept of fuel poverty. It also avoids reducing Westminster's mixed-income geography to area averages. A property can have moderate deterministic burden under the MSOA mean but still have a substantial probability of high burden for households in the local lower tail. Conversely, a property with high required cost may be less probabilistically exposed if local residual incomes are high. This distinction is central to the paper's question because price shocks will interact with both required demand and income sensitivity.

The residual gap between the modelled poverty proxy and official LSOA fuel-poverty rates is interpreted as a limitation of construct alignment rather than as a simple tuning failure. Income deprivation, required-energy burden, and official fuel poverty are related but non-equivalent measures. Under the current parameterisation, the MSOA mean term anchors $\bar{Y}_l$, while the LSOA lower-tail term strongly identifies $\sigma_l$; changing the lower-tail loss from squared error to pseudo-Huber L1 therefore did not materially change the fitted distributions. This is useful diagnostically because it cautions against using official fuel poverty as if it were the same object as income-tail exposure or future fuel price risk.

For robustness, the pipeline also reports a deterministic baseline burden ratio:
\begin{equation}
\widetilde{B}_{p,0}=\frac{C_{p,0}}{\bar{Y}_m},
\end{equation}
where $\bar{Y}_m$ is the MSOA after-housing-cost income point estimate. This measure is easier to interpret as a cost-to-income ratio, but it does not model within-LSOA income dispersion and is not calibrated to the local fuel-poverty comparison.

\subsection{EPC Enrichment and High-Risk Pathways}

The latest explanatory layer enriches the Westminster property master with EPC certificate records using UPRN-based linkage. UPRNs are standardised before both merging and QA checks, so that numeric and string representations of the same property identifier do not generate false unmatched cases. Where multiple EPC certificates are available for a property, the selected record is the latest \texttt{LODGEMENT\_DATETIME}; where this is unavailable, the latest \texttt{INSPECTION\_DATE} is used, and remaining ties fall back to the first available row. EPC fields are added with an \texttt{EPC\_} prefix, the original property master is not overwritten, and the row count and row order of the master are preserved.

The enrichment is followed by an explicit audit stage covering unmatched master records, duplicate UPRNs, added-field missingness, EPC rating distributions, lodgement-year distributions, and original versus corrected tenure groups. The EPC-derived tenure variable is treated as a proxy for tenure-related adaptive capacity rather than as observed legal tenure. The analysis version, \texttt{EPC\_TENURE\_GROUP\_FIXED}, classifies social rented records before private rented records to avoid misclassifying social rental descriptions as private rental cases.

The enriched stock is then transformed into a descriptive high-risk layer. Derived variables capture EPC performance, SAP-derived required energy cost, fabric weakness, electric heating exposure, gas connection, corrected tenure proxy, baseline required energy burden, and membership in modelled high-burden groups. The main current-condition groups are the top decile and top quintile of $P_{p,0}(0.10)$, denoted HR10 and HR20. Robustness groups use the top decile and top quintile of deterministic baseline required energy burden, denoted BR10 and BR20; an additional deterministic threshold group identifies properties with $\widetilde{B}_{p,0}>0.10$. These groups describe current modelled burden probability and required-burden intensity; they are not official fuel-poverty classifications and do not by themselves constitute the future-facing household fuel price risk measure.

The distinction between HR and BR groups is analytically useful rather than merely technical. HR groups identify properties whose required costs interact with the local LSOA income distribution to produce a high probability of crossing a 10 percent burden threshold. BR groups identify properties with high deterministic required burden relative to the MSOA point estimate of after-housing-cost income. Their overlap is therefore expected to be partial: the two approaches should identify a related high-burden tail, while also showing whether income-distribution assumptions change who is treated as exposed.

Finally, high-risk properties are assigned to mutually exclusive descriptive pathways using priority rules. The current typology distinguishes inefficient large owner-occupied gas houses, social rented estate or block risk, private rented weak-fabric gas flats, electric-heated inefficient flats or maisonettes, high-demand inefficient maisonettes, and a residual high-risk group. The pathway typology is not a causal model. Its purpose is to provide an interpretable bridge between property-level burden modelling and the later price-shock and intervention simulations, so that de-risking outcomes can be reported by mechanism and governance pathway rather than only for the aggregate stock.

The corresponding descriptive figures report pathway risk and contribution, EPC--floor-area sensitivity, tenure structure, electric exposure, spatial concentration, and HR--BR overlap (Figures~\ref{fig:hr-br-overlap}--\ref{fig:spatial-high-risk}). They are current-condition burden-exposure diagnostics rather than price-shock results.

\subsection{Price-Shock Risk Metrics}

The price-shock layer converts the baseline income--cost engine into forward-looking risk metrics. Let $s \in S$ denote a fuel-price scenario, including gas-price, electricity-price, combined-fuel, standing-charge, and policy-mitigation cases. Required cost under scenario $s$ is written as:
\begin{equation}
C_{p,s}=\pi^g_s Q^g_p+\pi^e_s Q^e_p+SC^g_s I^g_p+SC^e_s I^e_p-M_{p,s},
\end{equation}
where $Q^g_p$ and $Q^e_p$ are required gas and electricity quantities or archetype-imputed equivalents, $\pi^g_s$ and $\pi^e_s$ are scenario fuel prices, $SC^g_s$ and $SC^e_s$ are standing charges, $I^g_p$ and $I^e_p$ indicate fuel connections, and $M_{p,s}$ captures bill support or other mitigation. The scenario burden is:
\begin{equation}
B_{p,s}(Y_l)=\frac{C_{p,s}}{Y_l},
\end{equation}
and the burden uplift is:
\begin{equation}
\Delta B_{p,s}(Y_l)=B_{p,s}(Y_l)-B_{p,0}(Y_l).
\end{equation}

This cost equation is intentionally simple and explicit. It separates the components through which price risk reaches households: gas unit prices, electricity unit prices, fixed standing charges, fuel connections, and policy mitigation. That separation is necessary because different interventions act on different terms. Fabric retrofit reduces required quantities; heat pumps change the gas-electricity mix; PV and batteries reduce grid electricity imports under some profiles; tariff reform changes prices or standing charges; and public support enters as mitigation. A more compact total-bill equation would obscure these mechanisms and make it harder to explain why an intervention protects one group but not another.

The main risk indicators are defined as distributional rather than average outcomes. The scenario-specific threshold-crossing probability is:
\begin{equation}
P_{p,s}(\tau)=F_l\left(\frac{C_{p,s}}{\tau}\right).
\end{equation}
For scenario weights $\omega_s$, the expected threshold-crossing risk is:
\begin{equation}
R_{p}(\tau)=\sum_{s \in S}\omega_s P_{p,s}(\tau).
\end{equation}
An affordability gap-at-risk can be expressed as the expected shortfall above an acceptable burden threshold:
\begin{equation}
G_{p,s}(\tau)=\operatorname{E}_{Y_l}\left[\max\{C_{p,s}-\tau Y_l,0\}\right].
\end{equation}
Household Fuel Price-at-Risk is defined as a high quantile of the scenario burden distribution:
\begin{equation}
\operatorname{HFPaR}_{p,\alpha}
=\inf \left\{b:\sum_{s \in S}\omega_s \mathbf{1}\left(B_{p,s}\leq b\right)\geq \alpha \right\},
\end{equation}
where $\alpha$ may be set to a policy-relevant upper-tail level such as 0.90 or 0.95. These metrics keep fuel poverty, energy burden, and fuel price risk analytically separate: the first concerns current hardship, the second is a cost-to-income ratio, and the third measures future exposure and sensitivity under price shocks.

The three risk metrics have different policy uses. Threshold-crossing probabilities show how many property-income combinations are likely to move beyond an affordability boundary. Gap-at-risk reports the monetary depth of exposure above that boundary, which matters for bill support or social tariff design. Household Fuel Price-at-Risk focuses on the upper tail of the scenario burden distribution, which is the relevant object when the policy concern is protection against shocks rather than average annual affordability. Reporting all three avoids a false choice between headcount, severity, and tail-risk interpretations.

\subsection{Intervention Counterfactuals}

Domestic net zero interventions are represented as counterfactual changes to required demand, fuel structure, tariff exposure, and mitigation access. Let $k$ denote an intervention package, such as fabric retrofit, clean heat, PV, battery storage, smart tariffs, or an integrated package. The post-intervention scenario cost is:
\begin{equation}
C_{p,s,k}=
a_{p,k}\widetilde{C}_{p,s,k}+(1-a_{p,k})C_{p,s},
\end{equation}
where $a_{p,k}\in[0,1]$ represents technical, tenure, leasehold, heritage, roof, finance, or tariff feasibility. This term is central to the paper's interpretation: an intervention produces household risk protection only where access and market conditions allow the technical change to translate into lower exposure.

The feasibility parameter is not a decorative adjustment. It is the formal expression of the paper's claim that domestic net zero is conditional household protection. Without $a_{p,k}$, the model would implicitly assume that every property can receive every intervention and that all occupants can capture the resulting benefit. That assumption would be especially misleading in Westminster, where flats, leasehold blocks, private renting, social housing delivery, heritage constraints, roof rights, smart-meter access, and tariff eligibility all shape whether an engineering option becomes a household de-risking option.

The intervention-specific threshold risk and Fuel Price-at-Risk are:
\begin{equation}
P_{p,s,k}(\tau)=F_l\left(\frac{C_{p,s,k}}{\tau}\right),
\qquad
\operatorname{HFPaR}_{p,\alpha,k}
=Q_{\alpha}\left(\{B_{p,s,k}\}_{s \in S}\right).
\end{equation}
The de-risking dividend is the reduction in upper-tail burden risk:
\begin{equation}
D_{p,\alpha,k}=
\operatorname{HFPaR}_{p,\alpha,0}
-
\operatorname{HFPaR}_{p,\alpha,k}.
\end{equation}
Results will be stratified by income distribution, EPC rating, dwelling form, heating system, tenure, leasehold and heritage constraints, and LSOA/MSOA location. The empirical question is therefore not whether a low-carbon measure reduces emissions in engineering terms, but who receives measurable price-risk protection once required demand, tariffs, fuel-market pass-through, and local governance constraints are jointly considered.

The de-risking dividend is the bridge between the conceptual and empirical contributions. It does not ask whether an intervention is technically impressive on average. It asks whether the intervention reduces the upper-tail burden that households would face across price-shock scenarios. A positive dividend means that the package has converted at least part of the technical or policy change into lower household price-risk exposure. A small or uneven dividend indicates either that the intervention does not affect the relevant risk channel, or that feasibility, tariff access, price pass-through, or income sensitivity prevents households from capturing the protection.
